% !TeX root = ../../../../../main.tex
% chapters/sections/chapter3/subsections/households/stage2.tex

\subsection{第2段階（期をまたぐ問題）}
\label{sec:chapter3_households_stage2}
ここからは第2段階、すなわち期をまたぐ最適化問題を扱う。

\subsubsection{自国家計の第2段階}
\label{sec:chapter3_households_stage2_home}
第2段階（期をまたぐ最適化問題）を定式化するにあたり、第1段階である期内費用最小化の結果を用いて
期待生涯効用関数および予算制約式内の変数を集計変数へと整理する。
自国財消費指数\(c_t^{h\to H}\)、外国財消費指数\(c_t^{h\to F}\)および総消費指数\(c_t^{h\to W}\)の定義
\eqref{form:problem_setup_home_c_h_H_def}、
\eqref{form:problem_setup_home_c_h_F_def}、
\eqref{form:problem_setup_home_c_h_W_def}
より効用関数内の\(\log\)の引数の全体は総消費指数\(c_t^{h\to W}\)に一致するのであった。
\begin{equation}
\label{form:chapter3_households_stage2_home_log_argument_eq}
    \begin{aligned}
        &\frac{\left(\left[\sum_{h'\in H}(c_t^{h\to h'})^{\frac{\theta^H-1}{\theta^H}}\right]^{\frac{\theta^H}{\theta^H-1}}\right)^{\alpha^H}\left(\left[\sum_{f\in F}(c_t^{h\to f})^{\frac{\theta^F-1}{\theta^F}}\right]^{\frac{\theta^F}{\theta^F-1}}\right)^{1-\alpha^H}}{(\alpha^H)^{\alpha^H}(1-\alpha^H)^{1-\alpha^H}} \\
        &=\frac{(c_t^{h\to H})^{\alpha^H}(c_t^{h\to F})^{1-\alpha^H}}{(\alpha^H)^{\alpha^H}(1-\alpha^H)^{1-\alpha^H}}\quad\text{（\ref{form:chapter3_households_stage1a_home_c_h_H_definition_def}、\ref{form:chapter3_households_stage1b_home_c_h_W_definition_def}参照）} \\
        &=c_t^{h\to W}
    \end{aligned}
\end{equation}
次に予算制約式の支出側についても式
\eqref{form:chapter3_households_stage1b_cpi_descriptions_p_H_W_average_eq}
より、名目総支出\(\sum_{h'\in H}p_t^{h'}c_t^{h\to h'}+\sum_{f\in F}p_t^fc_t^{h\to f}\)は
消費者物価指数\(p_t^{H\to W}\)と総消費指数\(c_t^{h\to W}\)の積に等しくなるのであった。
\begin{equation}
\label{form:chapter3_households_stage2_home_p_H_W_average_eq}
    \begin{aligned}
        &\sum_{h'\in H}p_t^{h'}c_t^{h\to h'}+\sum_{f\in F}p_t^fc_t^{h\to f} \\
        &=p_t^{H\to W}c_t^{h\to W}
    \end{aligned}
\end{equation}
これらの関係を用いることで、自国家計の元となる最適化問題
\eqref{form:problem_setup_home_original_budget}
を二段階に分解したときの第二段階は、より単純な操作変数を用いた問題
\eqref{form:chapter3_households_problem_setup_home_stage2_max_def}
に書き換えることができる。
操作変数は総消費指数\(c_t^{h\to W}\)、自国コンティンジェント債券\(d_{t+1}^{h\to H}(j')\)、
および外国通貨建てリスクフリー債券\(b_{t+1}^{h\to F}\)であった。
（労働\(l_t^h\)は自身が設定する財の価格に応じて自動で決定されるため操作変数ではない）
この問題のラグランジアンは以下のようになる。
\begin{equation}
\label{form:chapter3_households_stage2_home_lagrangian_eq}
    \begin{aligned}
        \mathcal{L}^h=&\operatorname{E}_s\Biggl[\sum_{t=s}^{\infty}\left(\prod_{k=0}^{t-s-1}\beta_{s+k}^H\right)\Biggl\{\left(\log c_t^{h\to W}-\frac{\phi^H}{2}(l_t^h)^2\right) \\
        &\quad+\lambda_t^h\biggl(\Bigl(d_t^{h\to H}+(1+i_{t-1}^F)\frac{e_t^{/*}}{e_{t-1}^{/*}}b_t^{h\to F}+(1-\tau_t^H)p_t^hy_t^h+t_t^H\Bigr) \\
        &\quad-\Bigl(\sum_{j'\in J}q_{t,t+1}(j')d_{t+1}^{h\to H}(j')+b_{t+1}^{h\to F}+p_t^{H\to W}c_t^{h\to W}\Bigr)\biggr)\Biggr\}\Biggr]
    \end{aligned}
\end{equation}
まず総消費指数\(c_t^{h\to W}\)に関する偏微分から以下のFOCが得られる（式\eqref{xxx}TODO）。
\begin{equation}
\label{form:chapter3_households_stage2_home_lambda_h_p_H_W_c_h_W_eq}
    \lambda_t^h=\frac{1}{p_t^{H\to W}c_t^{h\to W}}
\end{equation}
自国コンティンジェント債券\(d_{t+1}^{h\to H}(j')\)に関する偏微分から得られたFOCを
\(j'\)について集計することで以下のオイラー方程式が導かれる（式\eqref{xxx}TODO）。
\begin{equation}
\label{form:chapter3_households_stage2_home_lambda_h_i_H_eq}
    \lambda_t^h=\beta_t^H(1+i_t^H)\operatorname{E}_t[\lambda_{t+1}^h]
\end{equation}
ここで\(i_t^H\)は国内の無リスク名目金利である。
外国通貨建てリスクフリー債券\(b_{t+1}^{h\to F}\)に関する偏微分から以下のオイラー方程式が得られる（式\eqref{xxx}TODO）。
\begin{equation}
\label{form:chapter3_households_stage2_home_lambda_h_i_F_eq}
    \lambda_t^he_t^{/*}=(1+i_t^F)\beta_t^H\operatorname{E}_t\left[\lambda_{t+1}^he_{t+1}^{/*}\right]
\end{equation}
ここで\(i_t^F\)は外国の無リスク名目金利である。

\subsubsection{外国家計の第2段階}
\label{sec:chapter3_households_stage2_foreign}
外国家計\(f\)についても自国家計と同様に第2段階（期をまたぐ最適化問題）を定式化するにあたり、
第1段階である期内費用最小化の結果を用いて期待生涯効用関数および予算制約式内の変数を集計変数へと整理する。
まず効用関数内の\(\log\)の引数は第1段階ステップAにおける各財消費指数の定義および
ステップBにおける総消費指数の定義に基づき、以下のように外国の総消費指数\(c_t^{f\to W}\)へと完全に集約されるのであった。
\begin{equation}
\label{form:chapter3_households_stage2_foreign_log_argument_eq}
    \begin{aligned}
        &\frac{\left(\left[\sum_{f'\in F}(c_t^{f\to f'})^{\frac{\theta^F-1}{\theta^F}}\right]^{\frac{\theta^F}{\theta^F-1}}\right)^{\alpha^F}\left(\left[\sum_{h'\in H}(c_t^{f\to h'})^{\frac{\theta^H-1}{\theta^H}}\right]^{\frac{\theta^H}{\theta^H-1}}\right)^{1-\alpha^F}}{(\alpha^F)^{\alpha^F}(1-\alpha^F)^{1-\alpha^F}} \\
        &=\frac{(c_t^{f\to F})^{\alpha^F}(c_t^{f\to H})^{1-\alpha^F}}{(\alpha^F)^{\alpha^F}(1-\alpha^F)^{1-\alpha^F}}\quad\text{（\ref{form:chapter3_households_stage1a_foreign_c_f_F_definition_def}、\ref{form:chapter3_households_stage1b_foreign_c_f_W_definition_def}参照）} \\
        &=c_t^{f\to W}
    \end{aligned}
\end{equation}
次に外国通貨建ての予算制約式の支出側についても式
\eqref{form:chapter3_households_stage1b_cpi_descriptions_p_F_W_star_average_eq}
より、名目総支出\(\sum_{f'\in F}p_t^{f'*}c_t^{f\to f'}+\sum_{h'\in H}p_t^{h'*}c_t^{f\to h'}\)は
外国の消費者物価指数\(p_t^{F\to W*}\)と総消費指数\(c_t^{f\to W}\)の積に等しくなるのであった。
\begin{equation}
\label{form:chapter3_households_stage2_foreign_p_F_W_star_average_eq}
    \begin{aligned}
        &\sum_{f'\in F}p_t^{f'*}c_t^{f\to f'}+\sum_{h'\in H}p_t^{h'*}c_t^{f\to h'} \\
        &=p_t^{F\to W*}c_t^{f\to W}
    \end{aligned}
\end{equation}
これらの関係を用いることで、外国家計の元となる最適化問題
\eqref{form:problem_setup_foreign_original_budget}
を二段階に分解したときの第二段階は、より単純な操作変数を用いた問題
\eqref{form:chapter3_households_problem_setup_foreign_stage2_max_def}
に書き換えることができる。
操作変数は総消費指数\(c_t^{f\to W}\)、外国コンティンジェント債券\(d_{t+1}^{f\to F*}(j')\)、
および自国通貨建てリスクフリー債券\(b_{t+1}^{f\to H*}\)であった。
（労働\(l_t^f\)は自身が設定する財の価格に応じて自動で決定されるため操作変数ではない）
この問題のラグランジアンは以下のようになる。
\begin{equation}
\label{form:chapter3_households_stage2_foreign_lagrangian_eq}
    \begin{aligned}
        \mathcal{L}^f=&\operatorname{E}_s\Biggl[\sum_{t=s}^{\infty}\left(\prod_{k=0}^{t-s-1}\beta_{s+k}^F\right)\Biggl\{\left(\log c_t^{f\to W}-\frac{\phi^F}{2}(l_t^f)^2\right) \\
        &\quad+\lambda_t^{f/*}\biggl(\Bigl(d_t^{f\to F*}+(1+i_{t-1}^H)\frac{e_{t-1}^{/*}}{e_t^{/*}}b_t^{f\to H*}+(1-\tau_t^F)p_t^{f*}y_t^f+t_t^{F*}\Bigr) \\
        &\quad-\Bigl(\sum_{j'\in J}q_{t,t+1}^*(j')d_{t+1}^{f\to F*}(j')+b_{t+1}^{f\to H*}+p_t^{F\to W*}c_t^{f\to W}\Bigr)\biggr)\Biggr\}\Biggr]
    \end{aligned}
\end{equation}
まず総消費指数\(c_t^{f\to W}\)に関する偏微分から以下のFOCが得られる（式\eqref{xxx}TODO）。
\begin{equation}
\label{form:chapter3_households_stage2_foreign_lambda_f_slash_star_p_F_W_star_c_f_W_eq}
    \lambda_t^{f/*}=\frac{1}{p_t^{F\to W*}c_t^{f\to W}}
\end{equation}
外国コンティンジェント債券\(d_{t+1}^{f\to F*}(j')\)に関する偏微分から得られたFOCを
\(j'\)について集計することで以下のオイラー方程式が導かれる（式\eqref{xxx}TODO）。
\begin{equation}
\label{form:chapter3_households_stage2_foreign_lambda_f_slash_star_i_F_eq}
    \lambda_t^{f/*}=\beta_t^F(1+i_t^F)\operatorname{E}_t[\lambda_{t+1}^{f/*}]
\end{equation}
ここで\(i_t^F\)は外国の無リスク名目金利である。
自国通貨建てリスクフリー債券\(b_{t+1}^{f\to H*}\)に関する偏微分から以下のオイラー方程式が得られる（式\eqref{xxx}TODO）。
\begin{equation}
\label{form:chapter3_households_stage2_foreign_lambda_f_slash_star_i_H_eq}
    \frac{\lambda_t^{f/*}}{e_t^{/*}}=(1+i_t^H)\beta_t^F\operatorname{E}_t\left[\frac{\lambda_{t+1}^{f/*}}{e_{t+1}^{/*}}\right]
\end{equation}
ここで\(i_t^H\)は自国の無リスク名目金利である。